\chapter{Méthodologie et Implémentation}
\label{chap:methode}

% --- CHAPEAU DU CHAPITRE ---
Dans ce chapitre, nous justifions nos choix d'implémentation afin de vérifier les trois hypothèses formulées précédemment. Ces choix résultent directement des constats de l'état de l'art. Nous détaillerons d'abord l'architecture technique retenue, puis le protocole rigoureux de création des stimuli linguistiques, pour enfin exposer les étapes d'analyse nous permettant de mesurer la dimensionnalité intrinsèque.

% --- SECTION 2.1 ---
\section{Justification des Choix Techniques}

Afin de mesurer l'évolution de la dimensionnalité intrinsèque (ID) au sein d'une architecture neuronale, nous avons sélectionné un ensemble d'outils spécifiques.

\subsection{Modèle de Langage : GPT-2 Small}
Nous avons choisi GPT-2 Small (via la librairie \textit{HuggingFace Transformers}) en raison de son architecture \textit{Decoder-only} causal qui traite l'information de manière unidirectionnelle (de gauche à droite, avec un masque causal). Ce choix est crucial pour étudier l'activité de \textit{ramping} (Hypothèse 1), car il simule la progression temporelle de la lecture humaine, contrairement aux modèles bidirectionnels comme BERT qui "voient" toute la phrase d'un coup. 

\subsection{Algorithme d'Estimation de l'ID : TwoNN}
Conformément à l'état de l'art établi par Facco et al. (2017) \cite{facco2017estimating}, nous utiliserons l'algorithme \textit{Two-Nearest Neighbors} (TwoNN). 

Son principal avantage est sa robustesse face aux variations de densité des données. Comme nous travaillons sur des variétés (\textit{manifolds}) locales (mot par mot), TwoNN permet d'estimer l'ID sans être biaisé par la courbure globale de l'espace des embeddings.
Le principe est le suivant : pour chaque point $i$ (représentation d'un token), on calcule les distances vers ses deux plus proches voisins, $r_{i,1}$ et $r_{i,2}$.

\subsection{Langage et Frameworks}
L’implémentation est réalisée en \textbf{Python}. Nous utilisons :
\begin{itemize}
    \item \textbf{PyTorch} pour l’extraction des \textit{hidden states} (états cachés) à chaque couche du modèle.
    \item \textbf{NumPy} pour la manipulation des données multidimensionnelles, le calcul matriciel et la gestion des tenseurs.
    \item \textbf{scikit-dimension} qui fournit une implémentation optimisée de l'algorithme TwoNN.
\end{itemize}

\subsection{Limite du modèle : absence de mécanisme de contrôle}
Contrairement au modèle MUC (voir Section~\ref{sec:muc}, page~\pageref{sec:muc}), l'architecture Transformer ne dispose pas de mécanisme de \og Contrôle \fg. Ce dernier, dans le cerveau humain, gère l'attention et oriente dynamiquement la sélection des informations pertinentes. GPT-2, en tant que modèle purement statistique, n'a pas cette capacité : face à du jabberwocky, il traite simplement une séquence improbable sans pouvoir adapter son attention en fonction du contexte sémantique. Les différences de dimensionnalité observées entre semantic et jabberwocky refléteront donc l'absence de sens, non une mobilisation active de ressources attentionnelles pour gérer le conflit.

% --- SECTION 2.2 ---
\section{Protocole de Génération des Stimuli}

La validation de nos hypothèses repose sur la comparaison de deux types de signaux. L’objectif est d’isoler l’impact du sens (la sémantique) en maintenant la structure grammaticale constante.

\subsection{Source des phrases}
Les phrases proviennent du corpus \textbf{WikiText-2} (Wikipedia, split train), en \textbf{anglais uniquement}. Ce choix est imposé par le fait que GPT-2 Small est entraîné exclusivement sur WebText (corpus anglophone). Toute phrase en français serait traitée comme une séquence incompréhensible, rendant la comparaison sémantique vs jabberwocky invalide.

Après normalisation (longueur 8--20 mots, ponctuation finale, pas de chiffres, URLs, ou parenthèses), un corpus de \textbf{$N = 500$ phrases} a été sélectionné pour le groupe sémantique. Ce nombre a été choisi pour deux raisons :
\begin{itemize}
    \item \textbf{Densité statistique pour TwoNN :} L'estimateur requiert suffisamment de points pour éviter les biais d'estimation. 500 phrases offrent une couverture appropriée.
    \item \textbf{Coût computationnel :} Extraction des hidden states et calcul de l'ID restent abordables sur CPU.
\end{itemize}

\subsection{Génération automatique du groupe "Jabberwocky"}
Pour tester l'Hypothèse 3, nous avons généré automatiquement une version "vide de sens" de chaque phrase sémantique. Ce procédé s'inspire des neurosciences cognitives pour présenter au modèle une séquence syntaxiquement structurée mais sémantiquement vide.

Le processus de transformation suit trois règles strictes (implémentées via spaCy) :
\begin{enumerate}
    \item \textbf{Préservation des mots-outils :} Tous les déterminants (DET), prépositions (ADP), conjonctions (CONJ), pronoms (PRON), auxiliaires (AUX) et particules (PART) sont conservés intacts.
    \item \textbf{Préservation des noms propres :} Pour réduire la dispersion BPE, les noms propres ne sont pas remplacés.
    \item \textbf{Substitution des mots porteurs :} Seuls les noms (NOUN), verbes (VERB), adjectifs (ADJ) et adverbes (ADV) sont remplacés par des pseudo-mots générés aléatoirement.
\end{enumerate}

Les pseudo-mots sont générés pour respecter les contraintes phonotactiques de l'anglais et sont globalement uniques (chaque pseudo-mot n'apparaît qu'une seule fois dans le corpus).

\textit{Exemple :} "The evolution of language is rarely taught" devient "The zumplation of language is zikly carting".

\subsection{Appariement et garanties}
Chaque phrase "Jabberwocky" possède exactement le même nombre de tokens et la même structure grammaticale que sa phrase "Sémantique" correspondante. Cette approche par paires garantit que toute différence de dimensionnalité observée reflète l'impact de l'absence de cohérence sémantique, dissociée de la longueur ou de la structure de surface.

\subsection{Enjeux de la Tokenisation (BPE)}
Avant d'être traitées, les phrases sont converties en unités numériques par l'algorithme \textit{Byte Pair Encoding} (BPE). Ce processus est crucial et induit un biais technique pour le Jabberwocky :
\begin{itemize}
    \item \textbf{Fragmentation :} Les mots du groupe "Sémantique" existent dans le vocabulaire et forment souvent un seul token. Les pseudo-mots (ex: \textit{vlipure}) sont inconnus et sont fragmentés en sous-unités (ex: \textit{vli + pure}).
    \item \textbf{Conséquence :} Cette fragmentation augmente mécaniquement le nombre de pas de temps. Nous devrons analyser l'évolution de l'ID de manière relative à la progression dans la phrase, car le modèle fournit un effort supplémentaire de reconstruction pour lier ces fragments.
\end{itemize}

% --- SECTION 2.3 ---
\section{Étapes de l'analyse}

L'analyse suit une séquence reproductible de scripts :

\subsection{Étape 0 : Validation de TwoNN (script \texttt{00\_validate\_twonn.py})}
Avant toute application au modèle, l'estimateur TwoNN est calibré et validé sur des jeux synthétiques de dimension intrinsèque connue. Cela permet de vérifier que l'implémentation produit des estimations robustes.

\subsection{Étape 1 : Génération des stimuli (script \texttt{01\_generate\_stimuli.py})}
Ce script extrait des phrases sémantiques du corpus WikiText-2, applique les filtres de nettoyage et génère les phrases Jabberwocky correspondantes en remplaçant les mots porteurs par des pseudo-mots respectant la phonotactique anglaise. Les paires de phrases sont sauvegardées au format JSON.

\subsection{Étape 2 : Extraction des représentations (script \texttt{02\_extract\_hidden\_states.py})}
Pour chaque phrase (sémantique + jabberwocky) :
\begin{itemize}
    \item Le modèle GPT-2 encode la phrase et produit des \textit{hidden states} à chaque couche (0--12) et position (token par token).
    \item Ces states sont alignés avec les mots de la phrase originale en gérant la fragmentation BPE (un jeton représentatif est sélectionné par décile).
    \item Les représentations vectorielles correspondantes sont sauvegardées au format NumPy compressé.
\end{itemize}

\subsection{Étape 3 : Calcul de l'ID (script \texttt{03\_compute\_id.py})}
Pour chaque condition (semantic, jabberwocky), couche, et décile de position :
\begin{itemize}
    \item On construit une matrice $\mathbf{X} \in \mathbb{R}^{N \times D}$ où chaque ligne est la représentation vectorielle d'une phrase à la position donnée.
    \item TwoNN estime l'ID de cette matrice locale, avec intervalles de confiance calculés par bootstrap sans remise.
    \item Les résultats sont testés statistiquement (Mann-Whitney avec correction Bonferroni, $\alpha = 0.05$ et régressions linéaires de ramping).
\end{itemize}

\subsection{Étape 4 : Visualisation (script \texttt{04\_visualize.py})}
Les profils d'ID par couche (Hunchback) et par décile (Ramping) sont tracés, ainsi que la distribution de la fragmentation BPE et la différence d'ID.

% --- CONCLUSION DU CHAPITRE (OBLIGATOIRE) ---
Dans ce chapitre, nous avons détaillé notre approche méthodologique, justifiant le choix des modèles (GPT-2, Llama) et de la métrique (TwoNN). Nous avons également défini un protocole strict de génération de stimuli (Sémantique vs Jabberwocky) pour isoler l'impact du sens sur la dimensionnalité. Ces outils et ces données sont maintenant prêts à être utilisés pour générer les résultats que nous analyserons dans le chapitre suivant.